% ---- Requires in 0main.tex: \usepackage{booktabs,siunitx,makecell,xcolor,listings}
% ---- Requires in 0main.tex: \lstset{...} with Python language settings

\section{Introduction}

Neural-PBIR~\cite{neuralpbir2024} is a multi-stage inverse-rendering pipeline
that recovers geometry, materials, and lighting from posed multi-view images.
The Physics-Based Inverse Rendering (PBIR) stage uses a differentiable path
tracer to refine materials and lighting via gradient-based optimization. Standard
PBIR uses a plain PathTracer for both the forward (loss-image) and backward
(gradient) passes, requiring many samples per pixel for low-variance gradients,
especially for specular materials.

PathReSTIR~\cite{restir2020} replaces the backward pass with spatial sample
resampling: candidate paths are reused across neighboring pixels to reduce
gradient variance at very low sample counts (1\,spp). This report evaluates
whether PathReSTIR can match standard PBIR quality while reducing the per-pixel
ray budget by half.

\section{Implementation}

The PathReSTIR integrator is implemented in C++ in the \texttt{psdr-jit}
library, extending the \texttt{PathIntegrator} base class with spatial reservoir
resampling. The \texttt{renderD()} method performs four passes:

\begin{enumerate}
    \item \textbf{Pass 0}: Primal rendering with spatial resampling (sample cache).
    \item \textbf{Pass 1}: Depth-0 ReSTIR gradient (backprop through first-bounce resampling).
    \item \textbf{Pass 2}: Depth-1+ MIS gradient (backprop through multi-importance sampling).
    \item \textbf{Pass 3}: Primary edge gradients (silhouette boundaries).
\end{enumerate}

The Python wrapper in \texttt{GREENFIELD/restir/run\_restir.py} defines a
\texttt{PathReSTIRRenderer} (\texttt{torch.autograd.Function}) that decouples
forward and backward: the forward pass calls PathTracer \texttt{renderC()} for
unbiased full-path loss images, while the backward pass calls PathReSTIR
\texttt{renderD()} for the spatial-resampling gradient.

\begin{lstlisting}[caption={PathReSTIRRenderer --- \texttt{torch.autograd.Function} wrapping PathTracer forward and PathReSTIR backward.},label={lst:renderer},float,floatplacement=t]
class PathReSTIRFunction(torch.autograd.Function):
    @staticmethod
    def forward(ctx, connector, scene, fwd_options, bwd_options,
                sensor_ids, fwd_integrator_id, bwd_integrator_id, *params):
        images = connector.renderC(scene, fwd_options,
                                   sensor_ids=sensor_ids,
                                   integrator_id=fwd_integrator_id)
        images = torch.nan_to_num(torch.stack(images, dim=0))
        ctx.connector = connector
        ctx.scene = scene
        ctx.bwd_options = bwd_options
        ctx.sensor_ids = sensor_ids
        ctx.bwd_integrator_id = bwd_integrator_id
        ctx.num_no_grads = 7
        return images

    @staticmethod
    def backward(ctx, grad_out):
        image_grads = [g.to("cuda") for g in grad_out]
        param_grads = ctx.connector.renderD(
            image_grads, ctx.scene, ctx.bwd_options,
            ctx.sensor_ids, ctx.bwd_integrator_id)
        return tuple([None] * ctx.num_no_grads + param_grads)
\end{lstlisting}

The \texttt{PSDRJITConnector} is patched via
\texttt{PSDRJITConnector.extensions} to register the \texttt{'path\_restir'}
integrator type and re-apply \texttt{use\_face\_normal} after mesh loading
(PathReSTIR requires this flag).

\section{Experimental Setup}

\subsection{ReSTIR Configuration}

The PBIR stage uses the same gin-configurable optimizer \texttt{opt.py} as the
standard pipeline, with these ReSTIR-specific settings:

\begin{itemize}
    \item Forward pass: PathTracer, 16\,spp
    \item Backward pass: PathReSTIR, 1\,spp, 8 candidates, radius 5
    \item 500 iterations, checkpoint every 25
    \item Microfacet basis BRDF, envmap lighting (latent-space)
\end{itemize}

\subsection{Dataset}

All 7 Stanford-ORB ``light'' evaluation scenes are used:
\texttt{teapot\_scene001}, \texttt{grogu\_scene002}, \texttt{gnome\_scene003},
\texttt{car\_scene004}, \texttt{pitcher\_scene005}, \texttt{blocks\_scene006},
\texttt{cactus\_scene007}. Each scene provides 200--300 multi-view HDR images
with ground-truth geometry, envmap lighting, and test/novel-view splits.

\subsection{Metrics}

Metrics are computed using the official Stanford-ORB evaluation harness
(\texttt{Stanford-ORB/scripts/test.py}):

\begin{itemize}
    \item \textbf{Novel-view synthesis}: PSNR (HDR and LDR), SSIM, LPIPS
    \item \textbf{Relighting}: PSNR (HDR and LDR), SSIM, LPIPS
    \item \textbf{Material (albedo)}: PSNR LDR, SSIM, LPIPS
    \item \textbf{Geometry}: Normal angle error, depth MSE
    \item \textbf{Shape}: Bidirectional Chamfer distance
\end{itemize}

All comparisons use the baseline Neural-PBIR published metrics (42-scene average).

\section{Bugs and Lessons Learned}

Several issues were discovered and fixed during the evaluation pipeline
development. These are documented here so future iterations avoid them.

\texttt{wandb.Json} \textbf{removed in wandb 0.16+.}
The initial code used \texttt{wandb.Json()} to serialize JSON results for
upload. This class was removed in wandb 0.16, causing an
\texttt{AttributeError} at the end of every benchmark run. Fixed by replacing
with \texttt{wandb.save()}.

\texttt{--scenes example} \textbf{routes only teapot.}
The Stanford-ORB test harness has a shorthand \texttt{example} that only
evaluates \texttt{teapot\_scene001}. When running multi-scene evaluations,
this silently produced incomplete results. Fixed by using \texttt{--scenes auto}
and filtering the eval input JSON to only include scenes with valid outputs.

\texttt{render\_geo} \textbf{not called by the evaluation pipeline.}
The \texttt{pipeline\_evaluate.py} script was missing the geometry rendering
stage (\texttt{render\_geo.py}), which meant depth and normal buffers were
never generated. The \texttt{eval\_preprocess.py} then asserted on missing
files. Fixed by adding a \texttt{render\_geo} stage before
\texttt{eval\_preprocess}.

\texttt{car\_scene004} \textbf{frame index offset.}
The scene's test frames start at index 64, not 0. The \texttt{render\_geo}
output filenames (\texttt{0064\_image.*}) didn't match the frame indices
expected by \texttt{eval\_preprocess.py}, causing assertion failures. This
was specific to scenes whose test split doesn't start at frame 0.

\textbf{Forward spp had surprisingly little impact.}
A three-way comparison (forward spp=4, 16, 64 with cand=8, 500 iters) on
\texttt{teapot\_scene001} showed less than 0.5\,dB difference between the
best (fwd=16) and worst (fwd=4) configurations. This suggests the backward
gradient quality from ReSTIR is the dominant factor, not the forward loss
image noise.

\section{Results}

\begin{table}[ht]
\centering
\caption{Stanford-ORB evaluation: PathReSTIR vs standard Neural-PBIR.}
\label{tab:results}
\footnotesize
\begin{tabular}{@{}lcccccc@{}}
\toprule
Scene &
  \makebox[1.0cm][c]{\shortstack{View\\PSNR HDR}} &
  \makebox[1.0cm][c]{\shortstack{View\\PSNR LDR}} &
  \makebox[1.0cm][c]{\shortstack{Light\\PSNR HDR}} &
  \makebox[1.0cm][c]{\shortstack{Light\\PSNR LDR}} &
  \makebox[1.0cm][c]{\shortstack{Mat.\\PSNR LDR}} &
  \makebox[1.0cm][c]{\shortstack{Shape\\$\times 10^{-5}$}} \\
\midrule
car\_scene004    & 22.34 & 27.90 & 23.68 & 29.73 & 40.42 & 3.79 \\
grogu\_scene002  & 20.92 & 26.42 & 22.28 & 28.34 & 39.32 & 43.0 \\
gnome\_scene003  & 22.93 & 28.72 & 18.87 & 25.08 & 32.94 & 3.63 \\
pitcher\_scene005& 24.18 & 30.75 & 20.90 & 27.15 & 37.21 & 48.0 \\
blocks\_scene006 & 23.92 & 28.65 & 24.12 & 29.06 & 36.74 & 6.80 \\
cactus\_scene007 & 24.83 & 32.23 & 26.81 & 35.08 & 46.80 & 2.08 \\
teapot\_scene001& 23.77 & 30.15 & 21.63 & 29.09 & 35.56 & 4.21 \\
\midrule
\textbf{ReSTIR avg (7)} & 23.3$\pm$1.2 & 29.3$\pm$1.8 & 22.6$\pm$2.4 & 29.1$\pm$2.8 & 38.4$\pm$4.1 & 16.0 \\
\midrule
\textbf{Std NPBIR (42)} & 28.8$\pm$4.6 & 36.8$\pm$3.5 & 26.0$\pm$3.4 & 33.3$\pm$2.9 & 38.7$\pm$4.0 & 21.5 \\
\bottomrule
\end{tabular}
\end{table}

Key observations:

\begin{itemize}
    \item \textbf{Shape and geometry} are essentially identical to the baseline.
    The Chamfer distance is actually slightly better for ReSTIR ($1.6\times10^{-4}$
    vs $2.1\times10^{-4}$), and normal-angle error is comparable
    ($0.076$\,rad vs $0.060$\,rad). These metrics are dominated by the NSR
    geometry stage, not the PBIR refinement.

    \item \textbf{Material} PSNR LDR is nearly identical ($38.4$ vs $38.7$),
    confirming that the BRDF and albedo optimization works well with ReSTIR
    gradients.

    \item \textbf{Visual metrics} (View and Light) are 5--6\,dB lower than the
    baseline. The forward pass at 16\,spp produces noisier loss images than the
    standard 64\,spp, which limits the optimizer's ability to match fine detail.

    \item \textbf{Speed}: The PBIR stage completes in approximately 100\,s per
    scene with ReSTIR versus approximately 210\,s for the standard PathTracer at
    the same 500 iterations --- a factor-of-2 reduction in ray budget (65
    rays/pixel vs 128).
\end{itemize}

\section{Next Steps}

\begin{itemize}
    \item Increase spatial candidates from 8 to 16 or 32 and re-evaluate.
    \item Increase forward-pass spp from 16 to 32 (still a ray saving vs 64).
    \item Tune spatial resampling radius and neighbor count per scene.
    \item Investigate why forward spp had such small impact (fwd=4 vs fwd=64).
    \item Apply temporal reuse (extending ReSTIR to ReSTIR DI) for the backward
    gradient pass.
\end{itemize}
